\section{Discussion}
\paragraph{What the extra rounds bought.} Two more rounds raised the best exact fitness from 0.722 to 0.764 and the test accuracy from 77.2\% to 79.4\%. Neither the test gain over the round-1 winner (12 vs.\ 8 discordant items) nor the one over Baphomet-1.0 (10 vs.\ 7) is significant. Baphomet-1.1 is the numerically best model we have measured, not a demonstrated improvement over Baphomet-1.0.

\paragraph{Winner's curse.} The winner is the maximum over many noisy measurements, so its measured value overstates its true quality. This happens twice here. First, the IRT estimate selected it for completion and overstated its exact fitness (0.782 vs.\ 0.764), as for the completed candidates on average ($+0.008$). Second, the exact 110-item fitness is itself a single sample at temperature 0.6 on a small item set, and the winner is the maximum of 52 exact values (40 round-1 candidates and 12 completions, the latter picked from 64 new candidates); part of its 0.042 lead over the round-1 winner on that set is likely selection noise. The 180-item test, which played no part in the selection, is the unbiased comparison, and there the lead shrinks to 2.2 points.

\paragraph{Efficiency.} Round 1 suggested that a capped generation fitness selects for short reasoning. Rounds 2--3 found a merge that is more accurate but somewhat longer than the round-1 winner (test: 1{,}834 vs.\ 1{,}498 tokens; GPQA truncation 28\% vs.\ 16\%), while still clearly shorter than Baphomet-1.0. At the 4{,}096-token fitness cap the two winners have similar lengths (1{,}149 vs.\ 1{,}096 tokens), so the difference appears mainly at the larger test cap.

\paragraph{IRT acceleration.} The 40-item estimator was accurate enough to steer the search (MAE 0.009, rank correlation 0.96 on completed candidates) at about 40\% of the tokens of a full evaluation. The completion rule kept the final choice on exact values, so the remaining IRT bias affects which candidates were completed, not the reported winner's fitness.

\paragraph{Fitness design matters.} Compared with the earlier log-likelihood search, which produced a merge that lost to Baphomet (2 vs.\ 17 discordant items on a 90-item subset) and broke tool calling, the generation-based fitness with a tool-use task produced merges that keep tool calling (BFCL 94--96\%) and match or exceed Baphomet-1.0. The runs differ in more than the fitness (merge method, donor count, search algorithm), so this comparison is suggestive only.

\section{Limitations}
One search with one seed; the 110-item fitness is one sample at temperature 0.6 and the winner was selected on it, so fitness-set comparisons with the donors are optimistic for Baphomet-1.1; candidates not completed (including the 3 stopped by racing) have only IRT estimates; the 180-item test is small (a 2--4-point difference is not significant); the donors and the uniform average were not run on the test; the fitness uses a shorter token cap than the test, and truncation explains part of every model's errors, so accuracies at a larger cap may differ; donor descriptions follow their model cards, which we did not verify.

\section{Conclusion}
An EvoGM-style search with a generation-based, token-capped fitness, continued for two rounds with an IRT-accelerated estimator under a further 6.53-GPU-hour run, produced Baphomet-1.1-4B. It is the best model on the fitness set (0.764 vs.\ 0.686 for the best donor and 0.603 for Baphomet-1.0, with selection bias in its favour) and numerically the best on the disjoint 180-item test (79.4\% vs.\ 77.8\% for Baphomet-1.0 and 70.0\% for Qwen3.5-4B), but not significantly better than Baphomet-1.0 ($p=0.63$). It uses 42\% fewer completion tokens than Baphomet-1.0. Next steps: a larger test set with several samples per item, accuracy--token budget curves (2k--16k), several search seeds, and test-set evaluation of the strongest donors.

\paragraph{Release.} Baphomet-1.1-4B is $\theta_0+\sum_i\lambda_{1.1,i}\tau_i$ with the coefficients above and base Qwen/Qwen3.5-4B. The released weights are reconstructed with the same fp8-quantized task vectors and the same merge code (\texttt{linmerge.py}) as the evaluated model, so they correspond to the weights tested here.

\paragraph{Compute.} Kaggle 2$\times$T4, vLLM 0.33.0. Round 1: 7.13 kernel hours (6.83\,h search, 0.3\,h test) plus 0.67\,h for a smoke test. Rounds 2--3: 6.53 kernel hours including setup, 64 subset evaluations, 12 completions and the 180-item test.

